% ============================================================================
\clearpage
\section*{v4.9.0-r63 Source-Only: The Frame Offset as a Ratio of Two Horizon Countings}
\addcontentsline{toc}{section}{v4.9.0-r63 Source-Only: The Frame Offset as a Ratio of Two Horizon Countings}
\label{sec:v490-r63-counting}
% ============================================================================

\paragraph{Purpose.}
Asked what \(\bsym\) measures, the author agreed to test a reading already offered in v4.8.0-r03 (F11): \(\gstring=\ln2/(\pi\sqrt3)\) is the value of the
Barbero--Immirzi parameter from a horizon counting with spin \(j=\tfrac12\) only (Ashtekar, Baez, Corichi, Krasnov), \(\gBI\) the value from the corrected counting with
all spins (Domagała--Lewandowski, Meissner).  r03 saw a conflict with the Lorentz reading.  This layer records the pre-registered check H47 (note
\texttt{llm/Vorab\_Notiz\_H47\_Zaehlregel\_beta\_sym.md}, commit \texttt{df0c788}), and, following the author's suggestions that the \(\tfrac12\) comes from the uncertainty
relation and that tachyons relax to a floor, two exact lemmas.  Earlier layers are not rewritten.

\noindent\textbf{Source-only status.}
This v4.9.0-r63 layer is source-only.  It does not build a PDF, does not assign a new DOI,
and does not authorize a public release.  The v4.9.0 Version DOI
\href{https://doi.org/10.5281/zenodo.23534954}{10.5281/zenodo.23534954} and the Concept DOI
\href{https://doi.org/10.5281/zenodo.19210728}{10.5281/zenodo.19210728} are unchanged by it.

\subsection*{v4.9.0-r63 results}
\label{subsec:v490-r63-results}

\paragraph{Theoremlet v4.9.0-r63.1 (H47).}
With the area spectrum \(a(j)=\sqrt{j(j+1)}\), let \(\gamma\) solve \(\sum_j g(j)\,e^{-2\pi\gamma a(j)}=1\).  Spin \(\tfrac12\) only, \(g=2\), gives \(\gamma_{\min}=\ln2/(\pi\sqrt3)\).
\begin{center}
\small
\begin{tabular}{@{}lcccc@{}}
\toprule
full counting & \(\gamma\) & \(\gamma/\gamma_{\min}\) & \(\beta\) & in \([0.545,0.570]\)\\
\midrule
\(g=2\) (Domagała--Lewandowski, Meissner) & \(0.2375330\) & \(1.864700\) & \(0.549406\) & yes\\
\(g=2j+1\) (Ghosh--Mitra) & \(0.2740669\) & \(2.151501\) & \(0.604468\) & no\\
\bottomrule
\end{tabular}
\end{center}
The main body uses the rounded input \(0.2375\) (\(\gL=1.86444\), \(\bsym=0.549348\)).  The window is the one of v4.9.0-r61.1(e).

\paragraph{Theoremlet v4.9.0-r63.2 (uncertainty saturation).}
In \(|j,m\rangle\) (\(\hbar=1\)) one has \(\Delta J_x\,\Delta J_y=\tfrac12\bigl(j(j+1)-m^2\bigr)\) and the Robertson bound \(\Delta J_x\,\Delta J_y\ge\tfrac12|m|\).  Equality holds
exactly for \(m=\pm j\), two states at every \(j\); at \(j=\tfrac12\) every state saturates and \(\Delta J_x=\tfrac12\).

\paragraph{Theoremlet v4.9.0-r63.3 (tachyon floor).}
A tachyon of invariant \(\mu\) has \(E=\mu/\sqrt{v^2-1}\), strictly decreasing for \(v>1\) with \(E\to0\) as \(v\to\infty\); the momentum tends to \(\mu\), the invariant does not
change.  The floor \(E=0\), \(p=\mu\) seen from a frame of velocity \(u\) has \(E'=-\gamma u\mu<0\), \(p'=\gamma\mu\), velocity \(-1/u\), and \(E'^2-p'^2=-\mu^2\).  In rapidities the
floor of a frame at \(\eta_X\) is \(\eta_X+i\pi/2\), with velocity \(\coth(\eta_X-\eta_Y)\) in a frame at \(\eta_Y\).

\paragraph{Theoremlet v4.9.0-r63.4 (where the floor sits).}
With M at \(+\eta\) and R at \(0\):
\begin{center}
\small
\begin{tabular}{@{}lccc@{}}
\toprule
floor of & velocity in R & velocity in M & \(\cosh\) of the real distance to M\\
\midrule
R & \(\infty\) & \(-1/\bsym\) & \(\cosh\eta=1.197\)\\
M & \(1/\bsym\) & \(\infty\) & \(1\)\\
mirror frame \(-\eta\) & \(-1/\bsym\) & \(-\coth2\eta\) & \(\cosh2\eta=\gL\)\\
\bottomrule
\end{tabular}
\end{center}
The tachyon of Convention~F, rapidity \(-\eta+i\pi/2\) (v4.9.0-r24.2), is the floor of the mirror frame.

\paragraph{Proof sketch.}
r63.1 by \texttt{llm/k102\_zaehlregeln.py} (output \texttt{llm/k102\_out.txt}).  r63.2: \(\langle J_x\rangle=\langle J_y\rangle=0\), \(\langle J_x^2\rangle=\langle J_y^2\rangle\), and
\(j(j+1)-m^2-|m|=(j-|m|)(j+|m|+1)\).  r63.3 by the boost formulas and \(\tanh(x+i\pi/2)=\coth x\).  r63.4 by r63.3.  The arithmetic of r63.2 and the boost statements of r63.3
are in \texttt{formal/QSFormal/Zaehlung.lean}; the operator values for \(j\le6\), the energy monotonicity and the table of r63.4 are in \texttt{llm/k103\_robertson\_tachyon.py}
(output \texttt{llm/k103\_out.txt}).  \(\square\)

\paragraph{Reading.}
\begin{enumerate}[leftmargin=2em]
\item Under this reading \(\bsym\) is not a measured quantity but an \(SU(2)\) number fixed by two counting rules.  The counting reading and the Lorentz reading of r03 agree
when the tachyon frame counts minimally and the matter frame fully: \(\gL=\cosh2\eta\) is the boost between these two opposite frames.
\item One principle gives both rules: a frame counts only states that saturate the uncertainty bound.  The tachyon frame takes the absolute minimum (\(j=\tfrac12\),
\(\Delta J_x=\tfrac12\)), the matter frame the minimum at each spin (\(m=\pm j\), weight \(2\)).  The Ghosh--Mitra counting also counts non-saturating states.
\item r63.3 gives a dynamical reason for the minimal counting: a tachyon that loses energy accelerates towards a single floor.
\item Only the floor of the mirror frame reproduces \(\gL\) as the real M--T distance.  A floor pinned to the matter frame gives distance \(0\), a floor in R gives
\(\cosh\eta\).  Why the floor sits in the mirror frame is open; the symmetric placement of Theorem~\ref{thm:symmetric-frame} is assumed, not derived.
\item r63.1 is not blind: the window was known, and the Ghosh--Mitra value had been estimated before the note.  OP1 now reads: derive from the frames why T counts only
\(j=\tfrac12\) and M all spins with weight \(2\).
\end{enumerate}

\subsection*{v4.9.0-r63 status in the three frames}
\label{subsec:v490-r63-frames}

\begin{center}
\small
\begin{tabular}{@{}>{\raggedright\arraybackslash}p{0.30\textwidth}>{\raggedright\arraybackslash}p{0.62\textwidth}@{}}
\toprule
result & M / T / P\\
\midrule
r63.1 & M counts all spins (weight \(2\)), T counts \(j=\tfrac12\); P carries no counting\\
r63.2 & frame-free (representation theory)\\
r63.3, r63.4 & kinematics; the T of Convention~F is the floor of the mirror frame\\
\bottomrule
\end{tabular}
\end{center}

\subsection*{v4.9.0-r63 register entries}
\label{subsec:v490-r63-register}

\begin{center}
\small
\begin{tabular}{@{}>{\raggedright\arraybackslash}p{0.34\textwidth}>{\raggedright\arraybackslash}p{0.58\textwidth}@{}}
\toprule
statement & status after v4.9.0-r63\\
\midrule
H47 (counting reading of \(\bsym\)) & DL/Meissner inside the lepton window, Ghosh--Mitra outside; not blind\\
uncertainty saturation & exact; weight \(2\) = states \(m=\pm j\)\\
tachyon floor & exact kinematics; floor of the mirror frame reproduces \(\gL\)\\
OP1 & open; reformulated as a question about two counting rules\\
\bottomrule
\end{tabular}
\end{center}

\subsection*{v4.9.0-r63 ledger and guard update}
\label{subsec:v490-r63-ledger-guard}

\[
R_{\rm cnt}^{v490\text{-}r63}=(1,1,0)
\]
for (counting reading consistent with the lepton window; both rules from uncertainty saturation; the assignment of the rules to the frames not derived).

The layer spends no guard:
\[
R_{\rm glob}^{(v490\text{-}r63,+)}=(0,1,1,1,1),
\qquad
\Gamma_{\rm guard}^{v490\text{-}r63}=(1,1,1,1,1),
\qquad
\beta_4=1/256\ \text{unspent}.
\]
It does not imply exact finite \(C5\) zero closure, exact all-readout terminality,
\(C6/T2\), a global photon-clock/RH theorem, or global all-readout closure.
